%%% ------------------------------------------------------------------
%%% JPC3 --- Improving the Czech Wordnet
%%% Project 3 (sensebank): the worked example.
%%%
%%% Every figure that comes from data is cited as \fact{key} from
%%% facts.tex.  To rebuild them:
%%%     cd sensebank
%%%     uv run pdt_wordnet_vs_wncze.py
%%%     uv run pdt_frequencies.py
%%%     uv run report/facts_tex.py
%%% then compile:
%%%     pdflatex report; bibtex report; pdflatex report; pdflatex report
%%% Figures quoted from the papers are typed, with their citation.
%%% ------------------------------------------------------------------

\documentclass[a4paper,12pt]{article}

\usepackage[margin=25mm]{geometry}
\usepackage[utf8]{inputenc}
\usepackage[T1]{fontenc}
\usepackage[english]{babel}

\usepackage[round]{natbib}
\bibliographystyle{plainnat}

\usepackage{booktabs}
\usepackage{graphicx}

\usepackage[svgnames]{xcolor}
\usepackage{hyperref,url}
\hypersetup{
    colorlinks=true,
    linkcolor=DarkBlue,
    urlcolor=DarkRed,
    citecolor=DarkGreen,
}

\usepackage{langsci-gb4e}
\let\eachwordone=\itshape
\providecommand{\noautomath}{}
\noautomath
\newcommand{\gm}[1]{\textsc{\MakeLowercase{#1}}}

\usepackage{langsci-avm}
\avmsetup{attributes=\scshape, values=\normalfont}

%%% Generated by facts_tex.py from the facts-*.json files:
%%% do not edit by hand.
\newcommand{\fact}[1]{%
  \ifcsname fact@#1\endcsname\csname fact@#1\endcsname
  \else\textbf{??}\PackageWarning{facts}{Undefined fact `#1'}\fi}
\newcommand{\factw}[1]{%
  \ifcsname factw@#1\endcsname\csname factw@#1\endcsname\else\fact{#1}\fi}

\expandafter\def\csname fact@cmp.absent.kept\endcsname{0}
\expandafter\def\csname factw@cmp.absent.kept\endcsname{zero}
\expandafter\def\csname fact@cmp.absent.lowconf\endcsname{120}
\expandafter\def\csname fact@cmp.absent.share\endcsname{0.9}
\expandafter\def\csname fact@cmp.absent.synsets\endcsname{383}
\expandafter\def\csname fact@cmp.absent.tagged\endcsname{905}
\expandafter\def\csname fact@cmp.absent.words\endcsname{560}
\expandafter\def\csname fact@cmp.all.kept\endcsname{26,351}
\expandafter\def\csname fact@cmp.all.lowconf\endcsname{303}
\expandafter\def\csname fact@cmp.all.share\endcsname{68.5}
\expandafter\def\csname fact@cmp.all.synsets\endcsname{19,494}
\expandafter\def\csname fact@cmp.all.tagged\endcsname{70,347}
\expandafter\def\csname fact@cmp.all.words\endcsname{26,351}
\expandafter\def\csname fact@cmp.brno.allkept\endcsname{2,757}
\expandafter\def\csname fact@cmp.brno.defs\endcsname{525}
\expandafter\def\csname fact@cmp.brno.defskept\endcsname{499}
\expandafter\def\csname fact@cmp.brno.examples\endcsname{44}
\expandafter\def\csname fact@cmp.brno.exampleskept\endcsname{44}
\expandafter\def\csname fact@cmp.brno.synsets\endcsname{3,133}
\expandafter\def\csname fact@cmp.brno.tagged\endcsname{7,934}
\expandafter\def\csname fact@cmp.defs.kept\endcsname{582}
\expandafter\def\csname fact@cmp.defs.total\endcsname{782}
\expandafter\def\csname fact@cmp.elsewhere.kept\endcsname{0}
\expandafter\def\csname factw@cmp.elsewhere.kept\endcsname{zero}
\expandafter\def\csname fact@cmp.elsewhere.lowconf\endcsname{88}
\expandafter\def\csname fact@cmp.elsewhere.share\endcsname{3.7}
\expandafter\def\csname fact@cmp.elsewhere.synsets\endcsname{1,024}
\expandafter\def\csname fact@cmp.elsewhere.tagged\endcsname{3,817}
\expandafter\def\csname fact@cmp.elsewhere.words\endcsname{1,481}
\expandafter\def\csname fact@cmp.examples.kept\endcsname{91}
\expandafter\def\csname fact@cmp.examples.total\endcsname{339}
\expandafter\def\csname fact@cmp.gone.kept\endcsname{0}
\expandafter\def\csname factw@cmp.gone.kept\endcsname{zero}
\expandafter\def\csname fact@cmp.gone.lowconf\endcsname{11}
\expandafter\def\csname fact@cmp.gone.share\endcsname{0.1}
\expandafter\def\csname fact@cmp.gone.synsets\endcsname{345}
\expandafter\def\csname fact@cmp.gone.tagged\endcsname{149}
\expandafter\def\csname fact@cmp.gone.words\endcsname{400}
\expandafter\def\csname fact@cmp.hyper.both\endcsname{19,441}
\expandafter\def\csname fact@cmp.hyper.kept\endcsname{13,155}
\expandafter\def\csname fact@cmp.hyper.orphans\endcsname{1,727}
\expandafter\def\csname fact@cmp.idcheck.overlap\endcsname{0}
\expandafter\def\csname factw@cmp.idcheck.overlap\endcsname{zero}
\expandafter\def\csname fact@cmp.idcheck.pwn15\endcsname{22,930}
\expandafter\def\csname fact@cmp.idcheck.same\endcsname{33}
\expandafter\def\csname fact@cmp.lemma.l0.now\endcsname{6}
\expandafter\def\csname factw@cmp.lemma.l0.now\endcsname{six}
\expandafter\def\csname fact@cmp.lemma.l0.pdt\endcsname{6}
\expandafter\def\csname factw@cmp.lemma.l0.pdt\endcsname{six}
\expandafter\def\csname fact@cmp.lemma.l0.reached\endcsname{5}
\expandafter\def\csname factw@cmp.lemma.l0.reached\endcsname{five}
\expandafter\def\csname fact@cmp.lemma.l0.tagged\endcsname{598}
\expandafter\def\csname fact@cmp.lemma.l1.now\endcsname{10}
\expandafter\def\csname fact@cmp.lemma.l1.pdt\endcsname{1}
\expandafter\def\csname factw@cmp.lemma.l1.pdt\endcsname{one}
\expandafter\def\csname fact@cmp.lemma.l1.reached\endcsname{1}
\expandafter\def\csname factw@cmp.lemma.l1.reached\endcsname{one}
\expandafter\def\csname fact@cmp.lemma.l1.tagged\endcsname{258}
\expandafter\def\csname fact@cmp.lemma.l10.now\endcsname{9}
\expandafter\def\csname factw@cmp.lemma.l10.now\endcsname{nine}
\expandafter\def\csname fact@cmp.lemma.l10.pdt\endcsname{4}
\expandafter\def\csname factw@cmp.lemma.l10.pdt\endcsname{four}
\expandafter\def\csname fact@cmp.lemma.l10.reached\endcsname{1}
\expandafter\def\csname factw@cmp.lemma.l10.reached\endcsname{one}
\expandafter\def\csname fact@cmp.lemma.l10.tagged\endcsname{2,803}
\expandafter\def\csname fact@cmp.lemma.l11.now\endcsname{1}
\expandafter\def\csname factw@cmp.lemma.l11.now\endcsname{one}
\expandafter\def\csname fact@cmp.lemma.l11.pdt\endcsname{6}
\expandafter\def\csname factw@cmp.lemma.l11.pdt\endcsname{six}
\expandafter\def\csname fact@cmp.lemma.l11.reached\endcsname{0}
\expandafter\def\csname factw@cmp.lemma.l11.reached\endcsname{zero}
\expandafter\def\csname fact@cmp.lemma.l11.tagged\endcsname{1,442}
\expandafter\def\csname fact@cmp.lemma.l12.now\endcsname{10}
\expandafter\def\csname fact@cmp.lemma.l12.pdt\endcsname{1}
\expandafter\def\csname factw@cmp.lemma.l12.pdt\endcsname{one}
\expandafter\def\csname fact@cmp.lemma.l12.reached\endcsname{1}
\expandafter\def\csname factw@cmp.lemma.l12.reached\endcsname{one}
\expandafter\def\csname fact@cmp.lemma.l12.tagged\endcsname{277}
\expandafter\def\csname fact@cmp.lemma.l13.now\endcsname{4}
\expandafter\def\csname factw@cmp.lemma.l13.now\endcsname{four}
\expandafter\def\csname fact@cmp.lemma.l13.pdt\endcsname{2}
\expandafter\def\csname factw@cmp.lemma.l13.pdt\endcsname{two}
\expandafter\def\csname fact@cmp.lemma.l13.reached\endcsname{1}
\expandafter\def\csname factw@cmp.lemma.l13.reached\endcsname{one}
\expandafter\def\csname fact@cmp.lemma.l13.tagged\endcsname{480}
\expandafter\def\csname fact@cmp.lemma.l14.now\endcsname{6}
\expandafter\def\csname factw@cmp.lemma.l14.now\endcsname{six}
\expandafter\def\csname fact@cmp.lemma.l14.pdt\endcsname{5}
\expandafter\def\csname factw@cmp.lemma.l14.pdt\endcsname{five}
\expandafter\def\csname fact@cmp.lemma.l14.reached\endcsname{4}
\expandafter\def\csname factw@cmp.lemma.l14.reached\endcsname{four}
\expandafter\def\csname fact@cmp.lemma.l14.tagged\endcsname{476}
\expandafter\def\csname fact@cmp.lemma.l15.now\endcsname{3}
\expandafter\def\csname factw@cmp.lemma.l15.now\endcsname{three}
\expandafter\def\csname fact@cmp.lemma.l15.pdt\endcsname{4}
\expandafter\def\csname factw@cmp.lemma.l15.pdt\endcsname{four}
\expandafter\def\csname fact@cmp.lemma.l15.reached\endcsname{1}
\expandafter\def\csname factw@cmp.lemma.l15.reached\endcsname{one}
\expandafter\def\csname fact@cmp.lemma.l15.tagged\endcsname{763}
\expandafter\def\csname fact@cmp.lemma.l16.now\endcsname{9}
\expandafter\def\csname factw@cmp.lemma.l16.now\endcsname{nine}
\expandafter\def\csname fact@cmp.lemma.l16.pdt\endcsname{4}
\expandafter\def\csname factw@cmp.lemma.l16.pdt\endcsname{four}
\expandafter\def\csname fact@cmp.lemma.l16.reached\endcsname{2}
\expandafter\def\csname factw@cmp.lemma.l16.reached\endcsname{two}
\expandafter\def\csname fact@cmp.lemma.l16.tagged\endcsname{217}
\expandafter\def\csname fact@cmp.lemma.l17.now\endcsname{4}
\expandafter\def\csname factw@cmp.lemma.l17.now\endcsname{four}
\expandafter\def\csname fact@cmp.lemma.l17.pdt\endcsname{4}
\expandafter\def\csname factw@cmp.lemma.l17.pdt\endcsname{four}
\expandafter\def\csname fact@cmp.lemma.l17.reached\endcsname{0}
\expandafter\def\csname factw@cmp.lemma.l17.reached\endcsname{zero}
\expandafter\def\csname fact@cmp.lemma.l17.tagged\endcsname{207}
\expandafter\def\csname fact@cmp.lemma.l18.now\endcsname{5}
\expandafter\def\csname factw@cmp.lemma.l18.now\endcsname{five}
\expandafter\def\csname fact@cmp.lemma.l18.pdt\endcsname{2}
\expandafter\def\csname factw@cmp.lemma.l18.pdt\endcsname{two}
\expandafter\def\csname fact@cmp.lemma.l18.reached\endcsname{2}
\expandafter\def\csname factw@cmp.lemma.l18.reached\endcsname{two}
\expandafter\def\csname fact@cmp.lemma.l18.tagged\endcsname{2,219}
\expandafter\def\csname fact@cmp.lemma.l19.now\endcsname{13}
\expandafter\def\csname fact@cmp.lemma.l19.pdt\endcsname{8}
\expandafter\def\csname factw@cmp.lemma.l19.pdt\endcsname{eight}
\expandafter\def\csname fact@cmp.lemma.l19.reached\endcsname{2}
\expandafter\def\csname factw@cmp.lemma.l19.reached\endcsname{two}
\expandafter\def\csname fact@cmp.lemma.l19.tagged\endcsname{1,002}
\expandafter\def\csname fact@cmp.lemma.l2.now\endcsname{6}
\expandafter\def\csname factw@cmp.lemma.l2.now\endcsname{six}
\expandafter\def\csname fact@cmp.lemma.l2.pdt\endcsname{3}
\expandafter\def\csname factw@cmp.lemma.l2.pdt\endcsname{three}
\expandafter\def\csname fact@cmp.lemma.l2.reached\endcsname{2}
\expandafter\def\csname factw@cmp.lemma.l2.reached\endcsname{two}
\expandafter\def\csname fact@cmp.lemma.l2.tagged\endcsname{688}
\expandafter\def\csname fact@cmp.lemma.l20.now\endcsname{6}
\expandafter\def\csname factw@cmp.lemma.l20.now\endcsname{six}
\expandafter\def\csname fact@cmp.lemma.l20.pdt\endcsname{4}
\expandafter\def\csname factw@cmp.lemma.l20.pdt\endcsname{four}
\expandafter\def\csname fact@cmp.lemma.l20.reached\endcsname{1}
\expandafter\def\csname factw@cmp.lemma.l20.reached\endcsname{one}
\expandafter\def\csname fact@cmp.lemma.l20.tagged\endcsname{422}
\expandafter\def\csname fact@cmp.lemma.l21.now\endcsname{4}
\expandafter\def\csname factw@cmp.lemma.l21.now\endcsname{four}
\expandafter\def\csname fact@cmp.lemma.l21.pdt\endcsname{1}
\expandafter\def\csname factw@cmp.lemma.l21.pdt\endcsname{one}
\expandafter\def\csname fact@cmp.lemma.l21.reached\endcsname{1}
\expandafter\def\csname factw@cmp.lemma.l21.reached\endcsname{one}
\expandafter\def\csname fact@cmp.lemma.l21.tagged\endcsname{216}
\expandafter\def\csname fact@cmp.lemma.l22.now\endcsname{10}
\expandafter\def\csname fact@cmp.lemma.l22.pdt\endcsname{4}
\expandafter\def\csname factw@cmp.lemma.l22.pdt\endcsname{four}
\expandafter\def\csname fact@cmp.lemma.l22.reached\endcsname{3}
\expandafter\def\csname factw@cmp.lemma.l22.reached\endcsname{three}
\expandafter\def\csname fact@cmp.lemma.l22.tagged\endcsname{364}
\expandafter\def\csname fact@cmp.lemma.l23.now\endcsname{5}
\expandafter\def\csname factw@cmp.lemma.l23.now\endcsname{five}
\expandafter\def\csname fact@cmp.lemma.l23.pdt\endcsname{2}
\expandafter\def\csname factw@cmp.lemma.l23.pdt\endcsname{two}
\expandafter\def\csname fact@cmp.lemma.l23.reached\endcsname{2}
\expandafter\def\csname factw@cmp.lemma.l23.reached\endcsname{two}
\expandafter\def\csname fact@cmp.lemma.l23.tagged\endcsname{422}
\expandafter\def\csname fact@cmp.lemma.l3.now\endcsname{4}
\expandafter\def\csname factw@cmp.lemma.l3.now\endcsname{four}
\expandafter\def\csname fact@cmp.lemma.l3.pdt\endcsname{2}
\expandafter\def\csname factw@cmp.lemma.l3.pdt\endcsname{two}
\expandafter\def\csname fact@cmp.lemma.l3.reached\endcsname{2}
\expandafter\def\csname factw@cmp.lemma.l3.reached\endcsname{two}
\expandafter\def\csname fact@cmp.lemma.l3.tagged\endcsname{791}
\expandafter\def\csname fact@cmp.lemma.l4.now\endcsname{15}
\expandafter\def\csname fact@cmp.lemma.l4.pdt\endcsname{10}
\expandafter\def\csname fact@cmp.lemma.l4.reached\endcsname{4}
\expandafter\def\csname factw@cmp.lemma.l4.reached\endcsname{four}
\expandafter\def\csname fact@cmp.lemma.l4.tagged\endcsname{515}
\expandafter\def\csname fact@cmp.lemma.l5.now\endcsname{4}
\expandafter\def\csname factw@cmp.lemma.l5.now\endcsname{four}
\expandafter\def\csname fact@cmp.lemma.l5.pdt\endcsname{2}
\expandafter\def\csname factw@cmp.lemma.l5.pdt\endcsname{two}
\expandafter\def\csname fact@cmp.lemma.l5.reached\endcsname{2}
\expandafter\def\csname factw@cmp.lemma.l5.reached\endcsname{two}
\expandafter\def\csname fact@cmp.lemma.l5.tagged\endcsname{375}
\expandafter\def\csname fact@cmp.lemma.l6.now\endcsname{7}
\expandafter\def\csname factw@cmp.lemma.l6.now\endcsname{seven}
\expandafter\def\csname fact@cmp.lemma.l6.pdt\endcsname{1}
\expandafter\def\csname factw@cmp.lemma.l6.pdt\endcsname{one}
\expandafter\def\csname fact@cmp.lemma.l6.reached\endcsname{0}
\expandafter\def\csname factw@cmp.lemma.l6.reached\endcsname{zero}
\expandafter\def\csname fact@cmp.lemma.l6.tagged\endcsname{469}
\expandafter\def\csname fact@cmp.lemma.l7.now\endcsname{12}
\expandafter\def\csname fact@cmp.lemma.l7.pdt\endcsname{7}
\expandafter\def\csname factw@cmp.lemma.l7.pdt\endcsname{seven}
\expandafter\def\csname fact@cmp.lemma.l7.reached\endcsname{0}
\expandafter\def\csname factw@cmp.lemma.l7.reached\endcsname{zero}
\expandafter\def\csname fact@cmp.lemma.l7.tagged\endcsname{179}
\expandafter\def\csname fact@cmp.lemma.l8.now\endcsname{3}
\expandafter\def\csname factw@cmp.lemma.l8.now\endcsname{three}
\expandafter\def\csname fact@cmp.lemma.l8.pdt\endcsname{2}
\expandafter\def\csname factw@cmp.lemma.l8.pdt\endcsname{two}
\expandafter\def\csname fact@cmp.lemma.l8.reached\endcsname{2}
\expandafter\def\csname factw@cmp.lemma.l8.reached\endcsname{two}
\expandafter\def\csname fact@cmp.lemma.l8.tagged\endcsname{268}
\expandafter\def\csname fact@cmp.lemma.l9.now\endcsname{15}
\expandafter\def\csname fact@cmp.lemma.l9.pdt\endcsname{5}
\expandafter\def\csname factw@cmp.lemma.l9.pdt\endcsname{five}
\expandafter\def\csname fact@cmp.lemma.l9.reached\endcsname{3}
\expandafter\def\csname factw@cmp.lemma.l9.reached\endcsname{three}
\expandafter\def\csname fact@cmp.lemma.l9.tagged\endcsname{468}
\expandafter\def\csname fact@cmp.newin.kept\endcsname{0}
\expandafter\def\csname factw@cmp.newin.kept\endcsname{zero}
\expandafter\def\csname fact@cmp.newin.lowconf\endcsname{0}
\expandafter\def\csname factw@cmp.newin.lowconf\endcsname{zero}
\expandafter\def\csname fact@cmp.newin.share\endcsname{3.3}
\expandafter\def\csname fact@cmp.newin.synsets\endcsname{93}
\expandafter\def\csname fact@cmp.newin.tagged\endcsname{3,386}
\expandafter\def\csname fact@cmp.newin.words\endcsname{113}
\expandafter\def\csname fact@cmp.newout.kept\endcsname{0}
\expandafter\def\csname factw@cmp.newout.kept\endcsname{zero}
\expandafter\def\csname fact@cmp.newout.lowconf\endcsname{0}
\expandafter\def\csname factw@cmp.newout.lowconf\endcsname{zero}
\expandafter\def\csname fact@cmp.newout.share\endcsname{0.3}
\expandafter\def\csname fact@cmp.newout.synsets\endcsname{71}
\expandafter\def\csname fact@cmp.newout.tagged\endcsname{319}
\expandafter\def\csname fact@cmp.newout.words\endcsname{88}
\expandafter\def\csname fact@cmp.newsyn.synsets\endcsname{164}
\expandafter\def\csname fact@cmp.newsyn.tagged\endcsname{3,705}
\expandafter\def\csname fact@cmp.noili.kept\endcsname{0}
\expandafter\def\csname factw@cmp.noili.kept\endcsname{zero}
\expandafter\def\csname fact@cmp.noili.lowconf\endcsname{0}
\expandafter\def\csname factw@cmp.noili.lowconf\endcsname{zero}
\expandafter\def\csname fact@cmp.noili.share\endcsname{0.1}
\expandafter\def\csname fact@cmp.noili.synsets\endcsname{90}
\expandafter\def\csname fact@cmp.noili.tagged\endcsname{64}
\expandafter\def\csname fact@cmp.noili.words\endcsname{112}
\expandafter\def\csname fact@cmp.some.kept\endcsname{2,748}
\expandafter\def\csname fact@cmp.some.lowconf\endcsname{41}
\expandafter\def\csname fact@cmp.some.share\endcsname{23.1}
\expandafter\def\csname fact@cmp.some.synsets\endcsname{1,594}
\expandafter\def\csname fact@cmp.some.tagged\endcsname{23,709}
\expandafter\def\csname fact@cmp.some.words\endcsname{5,054}
\expandafter\def\csname fact@cmp.total.kept\endcsname{29,099}
\expandafter\def\csname fact@cmp.total.synsets\endcsname{23,094}
\expandafter\def\csname fact@cmp.total.tagged\endcsname{102,696}
\expandafter\def\csname fact@cmp.total.words\endcsname{34,159}
\expandafter\def\csname fact@cmp.ufal.allkept\endcsname{68}
\expandafter\def\csname fact@cmp.ufal.defs\endcsname{155}
\expandafter\def\csname fact@cmp.ufal.defskept\endcsname{6}
\expandafter\def\csname factw@cmp.ufal.defskept\endcsname{six}
\expandafter\def\csname fact@cmp.ufal.examples\endcsname{242}
\expandafter\def\csname fact@cmp.ufal.exampleskept\endcsname{2}
\expandafter\def\csname factw@cmp.ufal.exampleskept\endcsname{two}
\expandafter\def\csname fact@cmp.ufal.synsets\endcsname{203}
\expandafter\def\csname fact@cmp.ufal.tagged\endcsname{17,430}
\expandafter\def\csname fact@freq.counted.share\endcsname{86.5}
\expandafter\def\csname fact@freq.exc.0\endcsname{166}
\expandafter\def\csname fact@freq.exc.1\endcsname{7,559}
\expandafter\def\csname fact@freq.exc.10\endcsname{1}
\expandafter\def\csname factw@freq.exc.10\endcsname{one}
\expandafter\def\csname fact@freq.exc.2\endcsname{761}
\expandafter\def\csname fact@freq.exc.4\endcsname{1,612}
\expandafter\def\csname fact@freq.exc.5\endcsname{259}
\expandafter\def\csname fact@freq.exc.6\endcsname{217}
\expandafter\def\csname fact@freq.exc.7\endcsname{2,386}
\expandafter\def\csname fact@freq.exc.8\endcsname{19}
\expandafter\def\csname fact@freq.fate.absent\endcsname{715}
\expandafter\def\csname fact@freq.fate.badtag\endcsname{0}
\expandafter\def\csname factw@freq.fate.badtag\endcsname{zero}
\expandafter\def\csname fact@freq.fate.counted\endcsname{88,854}
\expandafter\def\csname fact@freq.fate.exception\endcsname{12,980}
\expandafter\def\csname fact@freq.fate.lowconf\endcsname{1,640}
\expandafter\def\csname fact@freq.fate.newsynset\endcsname{3,704}
\expandafter\def\csname fact@freq.fate.nodoc\endcsname{3}
\expandafter\def\csname factw@freq.fate.nodoc\endcsname{three}
\expandafter\def\csname fact@freq.fate.noili\endcsname{64}
\expandafter\def\csname fact@freq.fate.notinsynset\endcsname{7,716}
\expandafter\def\csname fact@freq.fate.untagged\endcsname{33,189}
\expandafter\def\csname fact@freq.missing.pairs\endcsname{988}
\expandafter\def\csname fact@freq.missing.words\endcsname{7,716}
\expandafter\def\csname fact@freq.negated.0\endcsname{4.2}
\expandafter\def\csname fact@freq.negated.1\endcsname{0.0}
\expandafter\def\csname fact@freq.negated.10\endcsname{0.0}
\expandafter\def\csname fact@freq.negated.2\endcsname{0.0}
\expandafter\def\csname fact@freq.negated.4\endcsname{99.7}
\expandafter\def\csname fact@freq.negated.5\endcsname{0.4}
\expandafter\def\csname fact@freq.negated.6\endcsname{0.9}
\expandafter\def\csname fact@freq.negated.7\endcsname{0.0}
\expandafter\def\csname fact@freq.negated.8\endcsname{0.0}
\expandafter\def\csname fact@freq.negated.synset\endcsname{0.7}
\expandafter\def\csname fact@freq.negated.untagged\endcsname{2.2}
\expandafter\def\csname fact@freq.new.pairs\endcsname{53}
\expandafter\def\csname fact@freq.new.words\endcsname{3,704}
\expandafter\def\csname fact@freq.senses.both\endcsname{414}
\expandafter\def\csname fact@freq.senses.ntumc\endcsname{1,736}
\expandafter\def\csname fact@freq.senses.pdt\endcsname{7,080}
\expandafter\def\csname fact@freq.synsettags\endcsname{102,696}
\expandafter\def\csname fact@freq.words\endcsname{148,865}
\expandafter\def\csname fact@sheet.concept.judged\endcsname{80}
\expandafter\def\csname fact@sheet.concept.of\endcsname{80}
\expandafter\def\csname fact@sheet.items\endcsname{210}
\expandafter\def\csname fact@sheet.newsyn.judged\endcsname{53}
\expandafter\def\csname fact@sheet.newsyn.of\endcsname{53}
\expandafter\def\csname fact@sheet.noili.judged\endcsname{17}
\expandafter\def\csname fact@sheet.noili.of\endcsname{17}
\expandafter\def\csname fact@sheet.pairs\endcsname{21}
\expandafter\def\csname fact@sheet.perpair\endcsname{10}
\expandafter\def\csname fact@sheet.perstudent\endcsname{60}
\expandafter\def\csname fact@sheet.population\endcsname{1,138}
\expandafter\def\csname fact@sheet.students\endcsname{7}
\expandafter\def\csname factw@sheet.students\endcsname{seven}
\expandafter\def\csname fact@sheet.tagged\endcsname{4,817}
\expandafter\def\csname fact@sheet.word.judged\endcsname{60}
\expandafter\def\csname fact@sheet.word.of\endcsname{988}
\expandafter\def\csname fact@sheet.wordsample.once\endcsname{24}
\expandafter\def\csname fact@sheet.wordsample.tagged\endcsname{334}
\expandafter\def\csname fact@sheet.wordsample.ten\endcsname{11}


\title{Merging the PDT Sense Annotation into the Czech WordNet 2.2\\
  \large JPC3 --- Improving the Czech Wordnet}
\author{Francis Bond}
\date{2026-11-16 (draft: the hand validation is still to do)}

\begin{document}
\maketitle

\begin{abstract}
  Between 2004 and 2011 ÚFAL annotated \fact{freq.words} words of the
  Prague Dependency Treebank with synsets of their own copy of the Czech
  WordNet, and edited that copy as they went.  We ask how much of that work
  reached the current Czech WordNet, release 2.2, and use the annotation
  to give the release corpus frequencies.  Mapping the annotation
  wordnet's Princeton WordNet 1.5 ids through the Collaborative
  Interlingual Index, \fact{cmp.all.synsets} of its
  \fact{cmp.total.synsets} synsets reach a synset with all their words,
  and those synsets carry \fact{cmp.all.share}\% of the tags.  The edits made in Brno reached the release, but of the
  \fact{cmp.ufal.defs} definitions ÚFAL wrote, only
  \fact{cmp.ufal.defskept} did.
  The tags give counts for \fact{freq.senses.pdt} senses, from
  \fact{freq.counted.share}\% of the words tagged with a synset.
\end{abstract}

%%% ------------------------------------------------------------------
\section{Introduction}
%%% ------------------------------------------------------------------

In order to investigate meaning in the Prague Dependency Treebank,
there was a large scale annotation effort
\citep{hajic-etal2004,bejcek-etal2006,pdt-wsd2011} that assigned senses
to words int he treebank and also created/revised many synsets to make
appropriate meanings \cite{cwn19pdt2011}.  However, not all of this
information appears to have been kept in the main wordnet
\citep{rambousek-etal2017}.

For example, the PDT annotators tagged the Czech modal verb \textbf{\textit{moci}} \fact{cmp.lemma.l11.tagged}
times.  The two annotators disagreed on 1,435 of them, so ÚFAL rewrote
its senses \citep{bejcek-etal2006}: \factw{cmp.lemma.l11.pdt} new synsets
(possibility, permission, suitability, futility, estimate,
responsibility), each with a definition and examples such
as~(\ref{ex:moci}).  The Czech WordNet 2.2 has
\factw{cmp.lemma.l11.now} synset for \textit{moci}, with no definition or
example, and none of the \factw{cmp.lemma.l11.pdt} (Figure~\ref{fig:moci}).  Princeton wordnet does not include auxilliary verbs, so probably the Czech wordnet project is following this, but the NTU-MC wordnets have been adding them, and I think it would be good to have them in the wordnet.  In any case NAS-450034833 is synonymous with 00001740-a (i1) \textit{able} ``(usually followed by `to') having the necessary means or skill or know-how or authority to do something'': “able to swim”, “she was able to program her computer”, “we were at last able to buy a car”, “able to get a grant for the project”, even though it is a different POS.

FIXME:  wncze-NTUMC-80001788-v, with the words \textit{dát se, moci, moct} has no definition, no example and no ILI.  I thought it had a definition in the DB, I should check.  

\begin{figure}[ht]
  \centering\footnotesize
  \textbf{Czech WordNet 2.2}\par\medskip
  \avm{[ id & wncze-NTUMC-80001788-v \\
          lemmas & \textit{dát se}, \textit{moci}, \textit{moct} \\
          definition & --- \\
          source & NTU-MC \\
          example & --- \\
          frequency & 27 (NTU-MC) ]}
  \par\bigskip
  \textbf{PDT annotation wordnet}\par\medskip
  \begin{minipage}[t]{0.49\linewidth}
  \avm{[ id & NAS-450034833 \\
          lemmas & \textit{moci} \\
          definition & \textit{možnost, uskutečnitelnost} \\ & ``possibility, feasibility'' \\
          source & PDT, ÚFAL 2006 \\
          example & \textit{Každý se může zmýlit.} \\ & ``Anyone can make a mistake.'' \\
          frequency & 1,183 tags ]}
  \end{minipage}\hfill
  \begin{minipage}[t]{0.49\linewidth}
  \avm{[ id & NAS-753964023 \\
          lemmas & \textit{moci} \\
          definition & \textit{smět, dovolení} \\ & ``may, permission'' \\
          source & PDT, ÚFAL 2006 \\
          example & \textit{Může se tu kouřit?} \\ & ``Can one smoke here?'' \\
          frequency & 240 tags ]}
  \end{minipage}\par\medskip
  \begin{minipage}[t]{0.49\linewidth}
  \avm{[ id & NAS-1585073893 \\
          lemmas & \textit{moci} \\
          definition & \textit{odpovědnost} \\ & ``responsibility'' \\
          source & PDT, ÚFAL 2006 \\
          example & \textit{Já za to nemůžu} \\ & ``It's not my fault'' \\
          frequency & 9 tags ]}
  \end{minipage}\hfill
  \begin{minipage}[t]{0.49\linewidth}
  \avm{[ id & NAS-792515163 \\
          lemmas & \textit{moci} \\
          definition & \textit{vhodnost} \\ & ``suitability'' \\
          source & PDT, ÚFAL 2006 \\
          example & \textit{Mohl bys také pomoci.} \\ & ``You could help too.'' \\
          frequency & 6 tags ]}
  \end{minipage}\par\medskip
  \begin{minipage}[t]{0.49\linewidth}
  \avm{[ id & NAS-1493102466 \\
          lemmas & \textit{moci} \\
          definition & \textit{přibližný odhad (hovor.)} \\ & ``rough estimate (coll.)'' \\
          source & PDT, ÚFAL 2006 \\
          example & \textit{Může být tak pět hodin.} \\ & ``It may be about five o'clock.'' \\
          frequency & 4 tags ]}
  \end{minipage}\hfill
  \begin{minipage}[t]{0.49\linewidth}
  \avm{[ id & NAS-919753068 \\
          lemmas & \textit{moci} \\
          definition & \textit{marnost} \\ & ``futility'' \\
          source & PDT, ÚFAL 2006 \\
          example & \textit{Můžeš se třeba upracovat, \dots} \\ & ``Work as hard as you like, \dots'' \\
          frequency & 0 tags ]}
  \end{minipage}
  \caption{The senses of \textit{moci} in the Czech WordNet 2.2 (top) and
    in the PDT annotation wordnet (below), with the first example of each
    PDT synset.  PDT frequencies are the tags in the annotation; the
    release's is its NTU-MC count.  The release synset has no ILI.}
  \label{fig:moci}
\end{figure}

\begin{exe}
  \ex \label{ex:moci} Czech [ces] \\
      \gll Každý se může zmýlit. \\
           everyone \gm{refl} can.\gm{3sg} err.\gm{inf} \\
  \trans ``Anyone can make a mistake.'' \hfill \citep{cwn19pdt2011}
\end{exe}

In this paper we ask how much of the annotation's work reached the
current wordnet.  The tags are useful only as far as they link to a
wordnet people use, and for the release they are a second source of
sense frequencies, beside the NTU-MC corpus.

The project brief sets three goals.  Map the PDT tags to the current
synsets, aiming for at least 95\% automatically; use them to add
\emph{corpus frequencies} to the wordnet; and hand-categorise 500 of
the tags that do not map. Section~\ref{sec:previous} introduces
previous work, Section~\ref{sec:data} describes the data, and
Section~\ref{sec:method} what reached the release and the frequencies.
Section~\ref{sec:validation} reports the checks done so far and the
hand validation still to do.

%%% ------------------------------------------------------------------
\section{Previous work}
\label{sec:previous}
%%% ------------------------------------------------------------------

\citet{hajic-etal2004} describe the first year of the annotation.  The
Czech WordNet is the sense inventory; for each lemma the annotator sees
the synsets containing it and chooses one, preferring synsets where the
lemma is a single word, or one of eleven exceptions.  Two annotators tag
every word independently.  After a year they had 125,129 words, about
15\% of PDT, and report three findings.  Coverage was poor:
only 34\% of the words were in the wordnet, no adverbs at all, as the
version they used (1.2a) had none.  Agreement was low, 61.4\% on synsets
and 69.9\% counting exceptions, and did not fall as ambiguity rose.  The annotation also fed back into the wordnet: only 12\% of its synsets were
used, and very common senses of frequent words were missing.  They
present that feedback as the second purpose of the project, beside
training data for word sense disambiguation, and note that only SemCor
\citep{landes-etal1998} had done this before.

\citet{bejcek-etal2006} report the annotation at scale: 431,447 words read,
148,774 annotated, over 20\% of PDT 2.0.  Agreement counting exceptions
rose to 74.6\%.  They argue against letting annotators assign several
synsets: the exceptions for that case were agreed only 20 and 0 times.
Their main contribution is a correction procedure.  A third annotator
judges disagreements, and the 25 lemmas with at least 200 disagreements,
\textit{mít} with 2,853 and \textit{moci} with 1,435 among them, get
their senses rewritten in the wordnet, each with a definition and an
example, and are re-annotated.  Correcting 0.5\% of the problem lemmas
raised the words tagged with a single-word synset by 7.4 points and the
fully disambiguated sentences from 5,111 to 6,941, for 320 hours of a
corrector's time and 215 of a programmer's.  The data and the edited
wordnet were released in 2011 \citep{pdt-wsd2011,cwn19pdt2011}.  The released data agree: they have \fact{freq.words} annotated words,
against the 148,774 reported.

\citet{rambousek-etal2017} and \citet{rambousek-etal2018} tell the
wordnet's side.  The original Czech WordNet \citep{palasmrz2004} has
28,201 synsets linked to Princeton WordNet 2.0, and was published by ELRA
under a closed licence.  Later versions were never merged: a 2009 edit of
2,400 base-concept synsets, a 2011 extension from a bilingual dictionary
(40,621 synsets, 30\% of a checked sample with wrong synonyms), VerbaLex,
and 5,676 definitions written by students but kept in text files.  The
2017 roadmap is an open licence and links to CILI; the 2018 paper adds
derivational relations and a crowdsourcing tool for suggested fixes.
Neither mentions folding the PDT edits back in, though the 2018 paper
points to the PDT version on LINDAT as a ``slightly modified (corrected)''
copy of the original.  This project measures how much of them arrived.  The copies
had also drifted apart in their identifiers: the release's ids are
WordNet 2.0 offsets, the PDT copy's are 1.5 offsets.

%%% ------------------------------------------------------------------
\section{Data and resources}
\label{sec:data}
%%% ------------------------------------------------------------------

Table~\ref{tab:data} lists the five resources.  The three from LINDAT are
licensed CC BY-NC-SA, non-commercial, unlike the CC BY-SA of the Czech
WordNet release; Section~\ref{sec:release} returns to this.

\begin{table}[ht]
  \centering\small
  \begin{tabular}{@{}lll@{}}
    \toprule
    Resource & Used for & Licence \\
    \midrule
    Czech WordNet 1.9 PDT \citep{cwn19pdt2011} & the tagged synsets & CC BY-NC-SA 3.0 \\
    PDT sense annotation \citep{pdt-wsd2011} & the tags (stand-off) & CC BY-NC-SA 3.0 \\
    PDT 3.0 \citep{pdt30} & the lemmas & CC BY-NC-SA 3.0 \\
    CILI 1.5 map \citep{bond-etal2016cili} & offset to ILI & CC BY 4.0 \\
    Czech WordNet 2.2 \citep{pala-etal2026wncze} & the target & CC BY-SA 4.0 \\
    \bottomrule
  \end{tabular}
  \caption{Resources, what each is used for, and its licence.}
  \label{tab:data}
\end{table}

\subsection{The annotation wordnet}

  It has \fact{cmp.total.synsets}
synsets and \fact{cmp.total.words} words.  \fact{cmp.idcheck.pwn15} synsets
have Princeton WordNet ids, and \fact{cmp.newsyn.synsets} were added
without one.  The Princeton ids are \emph{1.5} offsets, not the 2.0
offsets the release uses: read as 2.0 ids, \fact{cmp.idcheck.same} of them
name a synset in the release, and \fact{cmp.idcheck.overlap} of those share a word
with it.  Read as 1.5 offsets through CILI they line up
(Section~\ref{sec:method}).  Edit stamps record who changed a synset and
when: Bejček, at ÚFAL in 2006, on \fact{cmp.ufal.synsets} synsets (the 203
``created or modified by ÚFAL'' of the LINDAT record), and the Brno team on
\fact{cmp.brno.synsets}.

\subsection{The annotation}

  One stand-off file per PDT document gives,
for each annotated word, its id and a tag.  The tag is a synset id without
leading zeros or part of speech (\texttt{1307561} for \texttt{01307561-v},
\textit{mít} ``have''), an id of a new synset, a small number for an
exception, or a hyphen for a word with no agreed tag.  The papers list the
exceptions but not their codes, so we inferred them (Table~\ref{tab:exc}).
Code~4 is the easiest to test: PDT lemmatises a negated word to its
positive form, so a word whose sense the wordnet has only in the positive
should be negated, and \fact{freq.negated.4}\% of code-4 words are, against
\fact{freq.negated.synset}\% of words tagged with a synset.  Code~8 has
\fact{freq.exc.8} words and code~10 has \factw{freq.exc.10}, close to the
20 and 0 agreed cases \citet{bejcek-etal2006} report for the exceptions
they number 8 and 10.

\begin{table}[ht]
  \centering\small
  \begin{tabular}{@{}rlrl@{}}
    \toprule
    Code & Exception & Words & Evidence; commonest lemmas \\
    \midrule
    1 & missing sense (guess)    & \fact{freq.exc.1} & none; \textit{kolo, hráč, kontakt} \\
    7 & proper name              & \fact{freq.exc.7} & lemmas; \textit{Profit} (a newspaper), \textit{Zeman} \\
    4 & missing negative sense   & \fact{freq.exc.4} & \fact{freq.negated.4}\% negated; \textit{chtít, vědět} \\
    2 & incorrect reflexivity    & \fact{freq.exc.2} & reflexive verbs; \textit{objevit (se), konat (se)} \\
    5 & incorrect lemma          & \fact{freq.exc.5} & lemmas; \textit{bank, haš} \\
    6 & figurative use           & \fact{freq.exc.6} & lemmas; \textit{dát, ruka, uzda} \\
    0 & other problem (guess)    & \fact{freq.exc.0} & none \\
    8 & unclear word meaning     & \fact{freq.exc.8} & count, as in \citet{bejcek-etal2006} \\
    10 & missing more general sense & \fact{freq.exc.10} & count, as in \citet{bejcek-etal2006} \\
    \bottomrule
  \end{tabular}
  \caption{Exception codes in the 2011 data, how each was identified, and
    its commonest lemmas.}
  \label{tab:exc}
\end{table}



%%% ------------------------------------------------------------------
\section{Method and results}
\label{sec:method}
%%% ------------------------------------------------------------------

\subsection{The bridge}

Each annotation synset with a Princeton id is mapped to an ILI through
the CILI table from WordNet~1.5 \citep{bond-etal2016cili}, which gives a
confidence for each link.  The release synset with that ILI is the
target.  We then compare the words, definitions, examples and hypernym
links of the two synsets, and count the tags on each annotation synset so
that every figure can be weighted by use
(\texttt{pdt\_wordnet\_vs\_wncze.py}).

\subsection{What reached the release}

Table~\ref{tab:status} gives the result.  \fact{cmp.all.synsets} synsets
reach a release synset with all their words, and they hold
\fact{cmp.all.share}\% of the tags; with partial matches it is
\fact{cmp.some.share}\% more.  The links below confidence~1 are where the
bridge breaks: \textit{dobrý} ``good'', \fact{cmp.lemma.l6.tagged} tags, maps
to \textit{výborný} ``excellent'' with confidence below 0.25.

\begin{table}[ht]
  \centering\small
  \begin{tabular}{@{}lrrrr@{}}
    \toprule
    What became of the synset & Synsets & Words & Tags & \% of tags \\
    \midrule
    all words kept                     & \fact{cmp.all.synsets} & \fact{cmp.all.words} & \fact{cmp.all.tagged} & \fact{cmp.all.share} \\
    some words kept                    & \fact{cmp.some.synsets} & \fact{cmp.some.words} & \fact{cmp.some.tagged} & \fact{cmp.some.share} \\
    no word kept, words elsewhere      & \fact{cmp.elsewhere.synsets} & \fact{cmp.elsewhere.words} & \fact{cmp.elsewhere.tagged} & \fact{cmp.elsewhere.share} \\
    no word kept, words gone           & \fact{cmp.gone.synsets} & \fact{cmp.gone.words} & \fact{cmp.gone.tagged} & \fact{cmp.gone.share} \\
    concept not in the release         & \fact{cmp.absent.synsets} & \fact{cmp.absent.words} & \fact{cmp.absent.tagged} & \fact{cmp.absent.share} \\
    no ILI for the 1.5 offset          & \fact{cmp.noili.synsets} & \fact{cmp.noili.words} & \fact{cmp.noili.tagged} & \fact{cmp.noili.share} \\
    new synset, words in one release synset & \fact{cmp.newin.synsets} & \fact{cmp.newin.words} & \fact{cmp.newin.tagged} & \fact{cmp.newin.share} \\
    new synset, not in the release     & \fact{cmp.newout.synsets} & \fact{cmp.newout.words} & \fact{cmp.newout.tagged} & \fact{cmp.newout.share} \\
    \midrule
    total                              & \fact{cmp.total.synsets} & \fact{cmp.total.words} & \fact{cmp.total.tagged} & \\
    \bottomrule
  \end{tabular}
  \caption{The annotation wordnet's synsets by what became of them in the
    Czech WordNet 2.2, weighted by the tags on them.  A new synset whose
    words share a release synset is weak evidence: for a single word it
    only means the word is there in some sense.}
  \label{tab:status}
\end{table}

Most of the editing done in Brno reached the release, and little of
ÚFAL's did.  Of
\fact{cmp.brno.defs} definitions in synsets Brno stamped,
\fact{cmp.brno.defskept} are in the release, and all
\fact{cmp.brno.examples} of their examples.  Of the \fact{cmp.ufal.defs}
definitions in synsets ÚFAL stamped, \fact{cmp.ufal.defskept} are, and
\fact{cmp.ufal.exampleskept} of \fact{cmp.ufal.examples} examples.  Those
\fact{cmp.ufal.synsets} synsets carry \fact{cmp.ufal.tagged} tags, the
most-used part of the annotation.  Of the hypernym links whose two ends
both reached the release, \fact{cmp.hyper.kept} of \fact{cmp.hyper.both}
are there; \fact{cmp.hyper.orphans} of the rest start from a release synset
with no hypernym at all, such as \textit{ampér} ``ampere''.

Table~\ref{tab:lemmas} follows the 25 corrected lemmas.  ``Reached''
counts the annotation senses whose ILI leads to a release synset that
contains the lemma.  For the two modal verbs, \textit{mít} and
\textit{moci}, almost nothing of the re-analysis arrived, and they are
among the most-tagged lemmas.

\begin{table}[ht]
  \centering\small
  \begin{tabular}{@{}llrrrr@{}}
    \toprule
    Lemma & Gloss & PDT senses & Tags & Reached & Release senses \\
    \midrule
    \textit{čas}    & time     & \fact{cmp.lemma.l0.pdt} & \fact{cmp.lemma.l0.tagged} & \fact{cmp.lemma.l0.reached} & \fact{cmp.lemma.l0.now} \\
    \textit{část}   & part     & \fact{cmp.lemma.l1.pdt} & \fact{cmp.lemma.l1.tagged} & \fact{cmp.lemma.l1.reached} & \fact{cmp.lemma.l1.now} \\
    \textit{cena}   & price    & \fact{cmp.lemma.l2.pdt} & \fact{cmp.lemma.l2.tagged} & \fact{cmp.lemma.l2.reached} & \fact{cmp.lemma.l2.now} \\
    \textit{člověk} & human    & \fact{cmp.lemma.l3.pdt} & \fact{cmp.lemma.l3.tagged} & \fact{cmp.lemma.l3.reached} & \fact{cmp.lemma.l3.now} \\
    \textit{dát}    & give     & \fact{cmp.lemma.l4.pdt} & \fact{cmp.lemma.l4.tagged} & \fact{cmp.lemma.l4.reached} & \fact{cmp.lemma.l4.now} \\
    \textit{den}    & day      & \fact{cmp.lemma.l5.pdt} & \fact{cmp.lemma.l5.tagged} & \fact{cmp.lemma.l5.reached} & \fact{cmp.lemma.l5.now} \\
    \textit{dobrý}  & good     & \fact{cmp.lemma.l6.pdt} & \fact{cmp.lemma.l6.tagged} & \fact{cmp.lemma.l6.reached} & \fact{cmp.lemma.l6.now} \\
    \textit{dostat} & get      & \fact{cmp.lemma.l7.pdt} & \fact{cmp.lemma.l7.tagged} & \fact{cmp.lemma.l7.reached} & \fact{cmp.lemma.l7.now} \\
    \textit{fax}    & fax      & \fact{cmp.lemma.l8.pdt} & \fact{cmp.lemma.l8.tagged} & \fact{cmp.lemma.l8.reached} & \fact{cmp.lemma.l8.now} \\
    \textit{místo}  & place    & \fact{cmp.lemma.l9.pdt} & \fact{cmp.lemma.l9.tagged} & \fact{cmp.lemma.l9.reached} & \fact{cmp.lemma.l9.now} \\
    \textit{mít}    & have     & \fact{cmp.lemma.l10.pdt} & \fact{cmp.lemma.l10.tagged} & \fact{cmp.lemma.l10.reached} & \fact{cmp.lemma.l10.now} \\
    \textit{moci}   & can      & \fact{cmp.lemma.l11.pdt} & \fact{cmp.lemma.l11.tagged} & \fact{cmp.lemma.l11.reached} & \fact{cmp.lemma.l11.now} \\
    \textit{návrh}  & proposal & \fact{cmp.lemma.l12.pdt} & \fact{cmp.lemma.l12.tagged} & \fact{cmp.lemma.l12.reached} & \fact{cmp.lemma.l12.now} \\
    \textit{podnik} & business & \fact{cmp.lemma.l13.pdt} & \fact{cmp.lemma.l13.tagged} & \fact{cmp.lemma.l13.reached} & \fact{cmp.lemma.l13.now} \\
    \textit{práce}  & work     & \fact{cmp.lemma.l14.pdt} & \fact{cmp.lemma.l14.tagged} & \fact{cmp.lemma.l14.reached} & \fact{cmp.lemma.l14.now} \\
    \textit{právo}  & law      & \fact{cmp.lemma.l15.pdt} & \fact{cmp.lemma.l15.tagged} & \fact{cmp.lemma.l15.reached} & \fact{cmp.lemma.l15.now} \\
    \textit{řada}   & row      & \fact{cmp.lemma.l16.pdt} & \fact{cmp.lemma.l16.tagged} & \fact{cmp.lemma.l16.reached} & \fact{cmp.lemma.l16.now} \\
    \textit{říkat}  & say      & \fact{cmp.lemma.l17.pdt} & \fact{cmp.lemma.l17.tagged} & \fact{cmp.lemma.l17.reached} & \fact{cmp.lemma.l17.now} \\
    \textit{rok}    & year     & \fact{cmp.lemma.l18.pdt} & \fact{cmp.lemma.l18.tagged} & \fact{cmp.lemma.l18.reached} & \fact{cmp.lemma.l18.now} \\
    \textit{stát}   & state; stand & \fact{cmp.lemma.l19.pdt} & \fact{cmp.lemma.l19.tagged} & \fact{cmp.lemma.l19.reached} & \fact{cmp.lemma.l19.now} \\
    \textit{svět}   & world    & \fact{cmp.lemma.l20.pdt} & \fact{cmp.lemma.l20.tagged} & \fact{cmp.lemma.l20.reached} & \fact{cmp.lemma.l20.now} \\
    \textit{systém} & system   & \fact{cmp.lemma.l21.pdt} & \fact{cmp.lemma.l21.tagged} & \fact{cmp.lemma.l21.reached} & \fact{cmp.lemma.l21.now} \\
    \textit{uvést}  & state    & \fact{cmp.lemma.l22.pdt} & \fact{cmp.lemma.l22.tagged} & \fact{cmp.lemma.l22.reached} & \fact{cmp.lemma.l22.now} \\
    \textit{vysoký} & high     & \fact{cmp.lemma.l23.pdt} & \fact{cmp.lemma.l23.tagged} & \fact{cmp.lemma.l23.reached} & \fact{cmp.lemma.l23.now} \\
    \bottomrule
  \end{tabular}
  \caption{The 25 lemmas ÚFAL re-analysed \citep[Table~4]{bejcek-etal2006}:
    senses in the annotation wordnet, tags on them, senses that reached a
    release synset with the lemma, and the lemma's senses in the release.
    \textit{stát} is the noun and the verb together.}
  \label{tab:lemmas}
\end{table}

\subsection{Corpus frequencies}

A tag is a count for a sense when its synset reaches a release synset by
a link of confidence~1, and that synset has the tagged word: the PDT lemma,
its reflexive form with \textit{se} or \textit{si}, or for a negated word
its form with \textit{ne-} (\texttt{pdt\_frequencies.py}).
Table~\ref{tab:fate} shows what became of every annotated word.
\fact{freq.counted.share}\% of the words tagged with a synset are counted,
giving counts for \fact{freq.senses.pdt} senses.  The release already has
counts from the NTU-MC Czech corpus for \fact{freq.senses.ntumc} senses;
\fact{freq.senses.both} senses have both.

\begin{table}[ht]
  \centering\small
  \begin{tabular}{@{}lr@{}}
    \toprule
    What became of the word & Words \\
    \midrule
    counted                                    & \fact{freq.fate.counted} \\
    no agreed tag                              & \fact{freq.fate.untagged} \\
    an exception (Table~\ref{tab:exc})         & \fact{freq.fate.exception} \\
    the word is not in the release synset      & \fact{freq.fate.notinsynset} \\
    a new synset, no Princeton id              & \fact{freq.fate.newsynset} \\
    CILI link below confidence 1               & \fact{freq.fate.lowconf} \\
    concept not in the release                 & \fact{freq.fate.absent} \\
    no ILI for the 1.5 offset                  & \fact{freq.fate.noili} \\
    the tag names no synset                    & \fact{freq.fate.badtag} \\
    the document is not in PDT 3.0             & \fact{freq.fate.nodoc} \\
    \midrule
    annotated words                            & \fact{freq.words} \\
    \bottomrule
  \end{tabular}
  \caption{The fate of each annotated word.}
  \label{tab:fate}
\end{table}

Table~\ref{tab:goals} sets the results against the goals.  The mapping
falls short of 95\%; most of the shortfall is words the release synset
lacks (\fact{freq.missing.words} words, \fact{freq.missing.pairs} pairs of
synset and lemma) and new synsets with no Princeton id, most of them
ÚFAL's (\fact{freq.new.words} words, \fact{freq.new.pairs} pairs).  Both lists are candidates for additions to
the release, which is the stretch goal.

\begin{table}[ht]
  \centering\small
  \begin{tabular}{@{}lll@{}}
    \toprule
    & Goal & Produced \\
    \midrule
    tags mapped automatically & $\geq$ 95\% & \fact{freq.counted.share}\% of synset tags \\
    senses with a frequency   & all mapped tags & \fact{freq.senses.pdt} \\
    residue hand-categorised  & 500 & \fact{sheet.items} senses, judged twice: TODO \\
    \bottomrule
  \end{tabular}
  \caption{What we set out to do, and what we did.}
  \label{tab:goals}
\end{table}

%%% ------------------------------------------------------------------
\section{Validation}
\label{sec:validation}
%%% ------------------------------------------------------------------

Four checks are done.  The id check and the negation test of
Section~\ref{sec:data} confirm the WordNet 1.5 ids and code~4.  The
definitions said to have reached the release are identical strings, not
near matches, in a random sample of six we read.  And the total of annotated words
agrees with \citet{bejcek-etal2006} to within 0.1\%.

The hand validation asks whether the senses the release lacks should be
added.  The author does not speak Czech, so \factw{sheet.students}
students of the course who do judge them, on sheets made by
\texttt{new\_sense\_sheets.py}.  The release lacks
\fact{sheet.population} PDT senses, a sense being a lemma with the
synset it was tagged with.  Table~\ref{tab:sheets} shows which of them
are judged: all those whose concept or synset is missing, and a random
\fact{sheet.word.judged} of the \fact{sheet.word.of} where only the word
is missing from the release synset.  Adding a word to a synset is easy,
so the sample serves only to measure how often it would be right; as it
is drawn uniformly, that share estimates the precision of adding all
\fact{sheet.word.of}.  The \fact{sheet.items} senses judged cover
\fact{sheet.tagged} tagged words.

For each sense a student sees up to three PDT sentences with the word
marked, the PDT synset with its definition and examples, its English
lemmas and gloss through the ILI, and the senses the release already has
for the lemma, each with its English gloss.  They choose one of four
answers: \emph{add}, a sense the release is missing; \emph{already in
the release}, giving which of its senses; \emph{wrong tag}, the
sentences do not use the word in the PDT sense; or \emph{unsure}.

The design is randomised, and every random step has its own seeded
stream, so the assignment is the same on every run.  Each sense goes to
one of the \fact{sheet.pairs} pairs of students, so every pair shares
\fact{sheet.perpair} senses and each student judges
\fact{sheet.perstudent}.  This gives Cohen's $\kappa$ for every pair, not
only overall, and shows up a student who reads the task differently.
The senses are dealt in order of frequency, in rounds of
\fact{sheet.pairs} to the pairs in a fresh random order, so every pair
judges both frequent and rare senses.  Item numbers and the order of each
sheet are random, and the sheets show neither how often a sense was
tagged nor why it is missing.

\begin{table}[ht]
  \centering\small
  \begin{tabular}{@{}lrrrrrrr@{}}
    \toprule
    Reason & Senses & Judged & Add & Already & Wrong & Unsure & $\kappa$ \\
    \midrule
    word not in the synset & \fact{sheet.word.of} & \fact{sheet.word.judged} & TODO & TODO & TODO & TODO & TODO \\
    concept not in the release & \fact{sheet.concept.of} & \fact{sheet.concept.judged} & TODO & TODO & TODO & TODO & TODO \\
    new synset, no Princeton id & \fact{sheet.newsyn.of} & \fact{sheet.newsyn.judged} & TODO & TODO & TODO & TODO & TODO \\
    no ILI for the 1.5 offset & \fact{sheet.noili.of} & \fact{sheet.noili.judged} & TODO & TODO & TODO & TODO & TODO \\
    \midrule
    total                                  & \fact{sheet.population} & \fact{sheet.items} & TODO & TODO & TODO & TODO & TODO \\
    \bottomrule
  \end{tabular}
  \caption{The PDT senses the release lacks, by why it lacks them; how
    many are judged; the students' answers; and agreement as Cohen's
    $\kappa$.}
  \label{tab:sheets}
\end{table}

TODO: when the sheets are back, fill in Table~\ref{tab:sheets}, report
$\kappa$ for each of the \fact{sheet.pairs} pairs, and resolve the
disagreements.  Then read the comments and the \emph{already} and
\emph{wrong tag} answers for why the release lacks each sense.
Table~\ref{tab:errors} lists the causes we have seen so far; count them
in the judged senses.

\begin{table}[ht]
  \centering\small
  \begin{tabular}{@{}lrl@{}}
    \toprule
    Cause & Count & Example \\
    \midrule
    wrong CILI link        & TODO & \textit{dobrý} ``good'' $\rightarrow$ \textit{výborný} ``excellent'' \\
    word dropped from the synset & TODO & \textit{práce} ``work'' not in \textit{robota, dřina, lopota} \\
    ÚFAL sense not in the release & TODO & \textit{moci} ``possibility'' \\
    senses split differently & TODO & \textit{mít}: \fact{cmp.lemma.l10.pdt} senses in PDT, \fact{cmp.lemma.l10.now} in the release \\
    hypernym missing in the release & TODO & \textit{ampér} has none \\
    \bottomrule
  \end{tabular}
  \caption{Why the release lacks a PDT sense: the causes seen so far, to
    count in the \fact{sheet.items} judged senses.}
  \label{tab:errors}
\end{table}

%%% ------------------------------------------------------------------
\section{Visualisation and explanation}
%%% ------------------------------------------------------------------

TODO: a flow diagram of the \fact{freq.words} annotated words, from the tag to
its fate (Table~\ref{tab:fate}), with the width of each band its number of
words.  It would show that, once a word has an agreed synset tag, the
mapping loses most words at two points: the missing words and the new
senses.

\begin{figure}[ht]
  \centering
  \fbox{\parbox[c][5cm][c]{.8\linewidth}{\centering TODO: flow of the annotated words}}
  \caption{Where the annotated words go.}
  \label{fig:flow}
\end{figure}

%%% ------------------------------------------------------------------
\section{Release}
\label{sec:release}
%%% ------------------------------------------------------------------

The frequencies are written as \fact{freq.senses.pdt} change records of the
Czech WordNet, one \texttt{set\_count} line per sense with corpus
\texttt{PDT}, ready for a pull request to the release repository.  The
missing words become \texttt{add\_sense} records once validated, and the
ÚFAL senses \texttt{add\_synset} records with their definitions and examples.

TODO: count how many of the annotation wordnet's definitions and examples
could be used in the release.  \fact{cmp.defs.total} definitions and
\fact{cmp.examples.total} examples are in it, and \fact{cmp.defs.kept} and
\fact{cmp.examples.kept} of them are in the release already; most of the
rest are ÚFAL's (\fact{cmp.ufal.defs} definitions and
\fact{cmp.ufal.examples} examples, \fact{cmp.ufal.defskept} and
\fact{cmp.ufal.exampleskept} in the release).  For each one missing, ask:
does its synset reach a release synset, or does it belong to a sense the
students judge worth adding; is it a definition of the concept or
an example filed as one (the second definition of \textit{rameno}
``branch'' is ``mrtvé rameno řeky''); and is an example a real sentence of the sense?
The answer is the number of \texttt{set\_definition} and
\texttt{add\_example} records to propose, under the licence below.

Before any of it goes in, the licence needs settling.  The annotation is
CC BY-NC-SA and the release CC BY-SA.  A count is a fact about the corpus,
but the senses and definitions ÚFAL wrote are their text, so we will ask
ÚFAL for permission, as Masaryk University gave it for the wordnet itself.

We will send ÚFAL the exception codes as we inferred them, for the
documentation of the data, and send CILI the WordNet 1.5 links the
validation finds wrong.

%%% ------------------------------------------------------------------
\section{Conclusion}
%%% ------------------------------------------------------------------

We have mapped the PDT sense annotation to the Czech WordNet 2.2.  The
synsets and their words came through: \fact{cmp.all.synsets} of
\fact{cmp.total.synsets} reach a release synset with all their words.  So
did the definitions written in Brno.  The senses, definitions and
examples ÚFAL wrote for the hard lemmas did not:
\fact{cmp.ufal.defskept} of \fact{cmp.ufal.defs} definitions are in the
release, and those senses are among the most tagged.  The tags give the
release counts for \fact{freq.senses.pdt} senses from a second corpus.
The scripts are in the \texttt{sensebank} directory of the JPC3
repository, and the counts go to the release as change records once ÚFAL
agrees to the licence.  Next, the students' judgements decide which of
the missing senses to add, starting with \textit{mít} and \textit{moci}.

%%% ------------------------------------------------------------------
\subsection*{Use of AI tools}
%%% ------------------------------------------------------------------

Claude Opus 5.5 (Anthropic) wrote the analysis scripts, inferred the
exception codes, and drafted this report, including the discussion of
the papers.  TODO: the author checks every result against the source data
and the papers before the report is final, and is responsible for the
content.

%%% ------------------------------------------------------------------

\bibliography{report}

\end{document}

%%% Local Variables:
%%% mode: latex
%%% TeX-master: t
%%% End:
